% ============================================================
% Abstract: The Information Budget of Graph Neural Networks
% ============================================================

% This file contains the abstract for the Information Budget paper
% To be included in main_information_budget.tex

% Abstract content (to be placed within \begin{abstract}...\end{abstract} in main file):
%
% Graph Neural Networks (GNNs) show inconsistent performance across datasets:
% sometimes improving over feature-only methods (MLPs) by 20%+, other times
% underperforming by similar margins. The conventional explanation attributes
% this to homophily, but datasets with identical homophily can show opposite
% outcomes. We discover that the critical factor is the Information Budget---
% the residual predictive capacity $(1 - \text{Acc}_{\text{MLP}})$ that features
% alone cannot explain.
%
% We establish three key principles: (1) The Conditional Enhancement Bound: Under
% fair comparison conditions (equivalent model capacity, identical features,
% standard aggregation), GNN improvement is upper-bounded by $(1 - \text{Acc}_{\text{MLP}})$,
% validated on 19 benchmark datasets (19/19 within bound) and 4 OGB datasets at
% million-node scale. (2) The Aggregation Damage Ratio (ADR): At mid-homophily
% ($h \approx 0.5$), GCN loses 18.9% of the information MLP captures. We prove
% via the Data Processing Inequality that concatenation-based architectures
% (GraphSAGE) cannot lose original feature information, explaining their
% 200$\times$ lower ADR. (3) Dual Heterophily Types: Type A (recoverable via
% 2-hop, H2GCN effective) vs Type B (irrecoverable, MLP optimal), resolving
% the Chameleon-Squirrel paradox.
%
% Our unified decision framework achieves 88.9% prediction accuracy on controlled
% CSBM experiments (32/36 configurations) and 77.8% on external validation (7/9
% datasets). This work shifts the focus from ``homophily determines GNN success''
% to ``feature sufficiency determines the ceiling, homophily determines the
% utilization''---providing practitioners with an actionable diagnostic: always
% train MLP first to reveal the budget.
